
\subsubsection{4.1 Dataset}

\textbf{Waterbirds} \citep{Sagawa2019Distributionally} is constructed by compositing bird images from CUB-200 \citep{Welinder2010CUB} onto backgrounds from Places \citep{Zhou2018Places}, creating a 95\% background-bird spurious correlation in the training set.

\begin{table*}[htbp]
\centering
\resizebox{\dimexpr 2\columnwidth+\columnsep\relax}{!}{%
\begin{tabular}{llllll}
\toprule
Split & Total & G0 (Landbird/Land) & G1 (Landbird/Water) & G2 (Waterbird/Land) & G3 (Waterbird/Water) \\
\midrule
Train & 4,795 & 3,498 (72.9\%) & 184 (3.8\%) & 56 (1.2\%) & 1,057 (22.1\%) \\
Val & 1,199 & --- & --- & --- & --- \\
Test & 5,794 & --- & --- & --- & --- \\
\bottomrule
\end{tabular}
}
\end{table*}

Group identity is encoded as G = y $\times$ 2 + place, where y $\in$ {0,1} is the bird class and place $\in$ {0,1} is the background type. Minority groups are G1 and G2, comprising approximately 5\% of training samples (240 of 4,795). Preprocessing: resize to $256\times 256$, center crop to $224\times 224$, ImageNet normalization (mean = [0.485, 0.456, 0.406], std = [0.229, 0.224, 0.225]).

\subsubsection{4.2 Model and Training}

\textbf{Model:} ResNet-50, pretrained on ImageNet-1K (ResNet50\_Weights.IMAGENET1K\_V1 API), with the final layer replaced by FC(2048 $\rightarrow$ 2).

\begin{table}[htbp]
\centering
\resizebox{\columnwidth}{!}{%
\begin{tabular}{ll}
\toprule
Hyperparameter & Value \\
\midrule
Optimizer & SGD \\
Learning rate & 0.001 \\
Momentum & 0.9 \\
Weight decay & $1\times 10^{-4}$ \\
Batch size & 128 \\
Training epochs & 10 \\
$\tilde{g}$ collection epochs & {1, 3, 5, 10} \\
Random seed & 42 \\
Hardware & NVIDIA H100 NVL \\
Framework & PyTorch 2.10+cu128, Python 3.10 \\
\bottomrule
\end{tabular}
}
\end{table}

All 4,795 training samples were evaluated at each collection epoch in eval() mode (BatchNorm statistics frozen). No data subsampling was performed.

\subsubsection{4.3 Evaluation Metrics}

\begin{table}[htbp]
\centering
\resizebox{\columnwidth}{!}{%
\begin{tabular}{lll}
\toprule
Metric & Definition & Target \\
\midrule
AUC & sklearn roc\_auc\_score; minority label = 1 for $G1\cup G2$ & > 0.70 \\
Ratio & mean $\tilde{g}(G1\cup G2)$ / mean $\tilde{g}(G0\cup G3)$ & $\geq$ 3.0x at T\_id=5; $\geq 1.2x$ at T\_id=10 \\
balance\_deviation & deviation of top-25\% subset from class-uniform distribution & $\leq$ 0.10 \\
h\_norm\_std\_ratio & std($\|h(x_{i})\|)$ / mean($\|h(x_{i})\|)$ per group & < 0.5 \\
\bottomrule
\end{tabular}
}
\end{table}

\subsubsection{4.4 Implementation}

The GradientNormAnalyzer class registers a forward hook on model.fc to capture FC input features. The outer-product decomposition $\tilde{g}_i = \|softmax(logit_{i}) -$ one\_hot($y_{i})\|$ is computed from saved features and model outputs. The codebase comprises approximately 786 lines across six source files (dataset.py: 122 lines; model.py: 89; train.py: 154; evaluate.py: 162; visualize.py: 259; run\_experiment.py: ~200). A test suite of 67 unit and integration tests was run; all 67 passed. Implementation planning compliance: 8/8 coding tasks completed (100\%).

